%%
% 第二章 相关工作（~7页）
%%

\chapter{相关工作}

本章从图神经网络与图基准、多模态属性图基准、多模态学术文档理解、学术文献图构建与自动化采集四个方面，系统综述与本文研究密切相关的工作，分析现有工作的局限性，明确本文的研究定位。

\section{图神经网络与图基准}

\subsection{图神经网络基础}

图神经网络是一类专门处理图结构数据的深度学习方法。给定一个图$G=(V, E)$，其中$V$为节点集合，$E$为边集合，每个节点$v_i \in V$具有特征向量$\mathbf{x}_i$。GNN的核心思想是通过消息传递（Message Passing）机制聚合邻居节点的信息来更新节点表示\cite{gilmer2017mpnn}。具体而言，在第$k$层迭代中，节点$v_i$的隐藏表示$\mathbf{h}_i^{(k)}$通过以下方式更新：
\begin{equation}
  \mathbf{h}_i^{(k)} = \text{UPDATE}^{(k)}\left(\mathbf{h}_i^{(k-1)}, \text{AGGREGATE}^{(k)}\left(\{\mathbf{h}_j^{(k-1)} : j \in \mathcal{N}(i)\}\right)\right)
  \label{eq:message_passing}
\end{equation}
其中$\mathcal{N}(i)$表示节点$v_i$的邻居集合，AGGREGATE函数负责聚合邻居信息，UPDATE函数负责结合自身信息和聚合后的邻居信息来更新节点表示。不同的AGGREGATE和UPDATE函数设计对应了不同的GNN模型。

\textbf{图卷积网络（GCN）。}Kipf和Welling\cite{kipf2017gcn}提出了图卷积网络，通过对图拉普拉斯矩阵的一阶切比雪夫多项式近似实现了高效的谱域图卷积。GCN的逐层传播规则为：
\begin{equation}
  \mathbf{H}^{(l+1)} = \sigma\left(\tilde{\mathbf{D}}^{-\frac{1}{2}}\tilde{\mathbf{A}}\tilde{\mathbf{D}}^{-\frac{1}{2}}\mathbf{H}^{(l)}\mathbf{W}^{(l)}\right)
  \label{eq:gcn}
\end{equation}
其中$\tilde{\mathbf{A}} = \mathbf{A} + \mathbf{I}$为添加自连接的邻接矩阵，$\tilde{\mathbf{D}}$为其度矩阵，$\mathbf{W}^{(l)}$为可学习参数，$\sigma$为非线性激活函数。GCN通过归一化的邻接矩阵实现了对称归一化的邻居聚合，其计算复杂度与图中边数呈线性关系，在大规模图上具有良好的可扩展性。GCN在引文网络分类等半监督节点分类任务上取得了当时的最佳性能，成为GNN领域最具影响力的工作之一。

\textbf{图注意力网络（GAT）。}Veli\v{c}kovi\'{c}等人\cite{velickovic2018gat}指出，GCN对所有邻居节点使用固定的聚合权重，无法自适应地区分不同邻居的重要性。为此，GAT引入了注意力机制，为每个邻居节点动态学习一个权重系数。注意力系数的计算方式为：
\begin{equation}
  \alpha_{ij} = \frac{\exp\left(\text{LeakyReLU}\left(\mathbf{a}^T[\mathbf{W}\mathbf{h}_i \| \mathbf{W}\mathbf{h}_j]\right)\right)}{\sum_{k \in \mathcal{N}(i)} \exp\left(\text{LeakyReLU}\left(\mathbf{a}^T[\mathbf{W}\mathbf{h}_i \| \mathbf{W}\mathbf{h}_k]\right)\right)}
  \label{eq:gat}
\end{equation}
其中$\mathbf{a}$为可学习的注意力向量，$\|$表示向量拼接。GAT还支持多头注意力机制，通过多组独立的注意力计算捕捉不同子空间的信息。注意力机制使GAT能够自适应地为重要邻居分配更高的权重，在存在噪声边的图上表现出更强的鲁棒性。

\textbf{GraphSAGE。}Hamilton等人\cite{hamilton2017graphsage}提出了GraphSAGE，解决了一个重要的实际问题——归纳式学习。与需要整图参与训练的GCN和GAT不同，GraphSAGE通过对邻居进行采样和聚合，能够对训练时未见过的节点生成表示。GraphSAGE支持多种聚合器（均值聚合、LSTM聚合、最大池化聚合），为每个节点采样固定数量的邻居，避免了计算开销随图规模增长的问题。

\textbf{图同构网络（GIN）。}Xu等人\cite{xu2019gin}从理论上分析了GNN的表达能力上界。他们证明，基于邻居聚合的GNN模型的表达能力不会超过Weisfeiler-Lehman（WL）图同构测试。基于这一理论分析，GIN采用求和聚合策略（而非均值或最大值聚合），使其在理论上达到最大的表达能力：
\begin{equation}
  \mathbf{h}_i^{(k)} = \text{MLP}^{(k)}\left((1 + \epsilon^{(k)}) \cdot \mathbf{h}_i^{(k-1)} + \sum_{j \in \mathcal{N}(i)} \mathbf{h}_j^{(k-1)}\right)
  \label{eq:gin}
\end{equation}
其中$\epsilon^{(k)}$为可学习的标量参数。

除上述经典模型外，Defferrard等人\cite{defferrard2016chebnet}提出的ChebNet通过切比雪夫多项式近似实现了快速的局部谱域滤波；Monti等人\cite{monti2017monet}提出的MoNet通过高斯混合模型定义了统一的空域聚合框架。Gilmer等人\cite{gilmer2017mpnn}提出了消息传递神经网络（MPNN）框架，将上述模型统一为消息传递的范式，为后续研究提供了清晰的理论框架。这些模型的不断发展推动了GNN在节点分类、图分类、链接预测等任务上的性能持续提升\cite{zhou2020gnn_survey, wu2021gnn_survey}。

\subsection{图基准数据集的发展}

高质量的基准数据集是推动GNN研究发展的关键基础设施。在图机器学习领域，基准数据集的发展经历了从经典小规模数据集到大规模标准化基准的演进过程。

\textbf{经典引文网络数据集。}Sen等人\cite{sen2008collective}系统整理了Cora、CiteSeer和PubMed三个引文网络数据集，长期以来被作为GNN模型评估的标准基准。Cora数据集包含2,708篇机器学习论文和5,429条引用边，节点属性为1,433维的词袋向量，标签为7个研究领域。CiteSeer数据集包含3,327篇论文和4,732条引用边，节点属性为3,703维词袋向量。PubMed数据集规模稍大，包含19,717篇糖尿病相关论文和44,338条引用边。这三个数据集在GNN研究的早期阶段发挥了重要作用，GCN\cite{kipf2017gcn}、GAT\cite{velickovic2018gat}等经典模型均在这些数据集上验证了其有效性。然而，这些数据集规模偏小，节点属性仅为低维词袋表示，难以充分评估模型在大规模图和丰富属性场景下的性能。

\textbf{开放图基准（OGB）。}Hu等人\cite{hu2020ogb}提出的OGB是图机器学习领域最重要的基准之一。OGB提供了涵盖节点属性预测、链接属性预测和图属性预测三大类任务的标准化评估框架。在节点分类任务中，ogbn-arxiv数据集包含约169,343篇arXiv论文节点和1,166,243条引用边，覆盖40个学术类别。ogbn-arxiv将论文的标题和摘要通过skip-gram模型编码为128维向量作为节点特征。ogbn-mag数据集则是一个异构学术图，包含约1.9M个节点（论文、作者、机构、研究领域四种类型）和约21M条边。OGB的重要贡献在于提供了统一的评估协议、自动化的性能排行榜和标准化的数据加载接口，极大地促进了图学习研究的可复现性和可比性。

然而，OGB在节点属性的丰富度上存在明显不足。ogbn-arxiv的128维skip-gram向量丢失了论文原始文本中的大量语义信息（如上下文语义、专业术语等），更无法体现论文中的视觉信息（实验图表、数学公式等）。这种信息损失限制了多模态图学习方法的研究和发展。

\textbf{文本属性图基准。}为了弥补OGB在文本信息方面的不足，Yan等人\cite{yan2023cstag}提出了CS-TAG基准，将论文的原始标题和摘要作为文本属性图（Text-Attributed Graph, TAG）的节点特征。CS-TAG包含Coauthor-CS和Coauthor-Phy两个数据集，支持在原始文本上评估不同的文本编码策略对图学习性能的影响。CS-TAG的实验表明，使用预训练语言模型（如SciBERT\cite{beltagy2019scibert}）编码文本可以显著提升节点分类性能。然而，CS-TAG仍局限于文本单模态，未引入图像和公式等其他模态信息。

Li等人\cite{li2024tegdb}提出了TEG-DB基准，首次将文本属性扩展到图的边上，支持边文本图学习任务。TEG-DB包含多个数据集，涵盖了论文引用（边文本为引用上下文）、社交互动（边文本为对话内容）等多种场景。然而，TEG-DB同样未引入多模态节点属性。

综上所述，现有图基准数据集在节点属性的丰富度上仍存在不足，特别是缺乏多模态节点属性（如图像、公式等）。表\ref{tab:graph_benchmarks}对主要图基准数据集的特性进行了系统对比。

\begin{table}[htbp]
  \centering
  \caption{主要图基准数据集特性对比}
  \label{tab:graph_benchmarks}
  \begin{tabular}{lccccc}
    \toprule
    数据集 & 节点数 & 节点属性 & 模态类型 & 图结构 & 规模 \\
    \midrule
    Cora\cite{sen2008collective} & 2,708 & 词袋(1433d) & 文本 & 引用网络 & 小 \\
    CiteSeer\cite{sen2008collective} & 3,327 & 词袋(3703d) & 文本 & 引用网络 & 小 \\
    ogbn-arxiv\cite{hu2020ogb} & 169K & skip-gram(128d) & 文本向量 & 引用网络 & 大 \\
    CS-TAG\cite{yan2023cstag} & 18K & 原始文本 & 文本 & 合作网络 & 中 \\
    TEG-DB\cite{li2024tegdb} & 多 & 节点+边文本 & 文本 & 多种类型 & 中 \\
    \bottomrule
  \end{tabular}
\end{table}

\section{多模态属性图基准}

多模态属性图（Multimodal Attributed Graph, MAG）是图结构与多模态节点属性的结合。与仅使用单一模态节点特征的传统图学习不同，MAG学习旨在利用多模态信息提升图学习任务的性能。Baltru\v{s}aitis等人\cite{baltruvsaitis2019multimodal}系统综述了多模态学习的发展历程和技术挑战，指出多模态融合、对齐和表示学习是多模态学习的三个核心问题。在图学习领域，如何有效地融合多模态节点属性与图结构信息，是一个新兴且具有挑战性的研究方向。

\subsection{MAGB：首个多模态属性图基准}

Yan等人\cite{yan2025magb}提出的MAGB是多模态属性图学习领域的开创性工作。MAGB首次系统地定义了MAG学习范式，并构建了包含5个数据集的基准，覆盖电商领域（Amazon Movies、Amazon Toys、Amazon Grocery）和社交领域（Reddit）。每个数据集中的节点都具有文本描述和产品/帖子图像两种模态属性。在文本编码方面，MAGB使用Sentence-BERT\cite{reimers2019sbert}将文本编码为384维向量；在图像编码方面，使用CLIP\cite{radford2021clip}将图像编码为512维向量。多模态融合采用简单拼接策略，将文本和图像向量拼接为896维的联合表示。

MAGB系统评估了16个基线模型在节点分类和链接预测任务上的表现，揭示了若干重要发现：（1）不同领域对模态的偏好存在显著差异——社交网络（Reddit）中视觉特征的贡献大于文本特征，而电商领域（Amazon）则更偏好文本特征；（2）简单拼接多模态特征存在模态不平衡问题，图像模态的高维度可能在某些场景下引入噪声；（3）图结构信息在多模态设置下仍能提供额外的增益，但这种增益因数据集而异。

然而，MAGB存在三个明显的局限：第一，领域覆盖局限于电商和社交领域，未涉及学术论文这一天然的多模态属性图场景；第二，模态类型仅限于文本和图像，不包含公式这一学术论文特有的模态；第三，部分数据集（如Amazon系列）不包含天然的图结构，其图结构是基于用户共同购买或共同评论关系构建的，与引用网络等学术图结构有本质差异。

\subsection{OpenMAG：大规模综合基准}

Wan等人\cite{wan2026openmag}提出了OpenMAG，在规模和评估维度上大幅扩展了MAG基准。OpenMAG包含19个数据集，覆盖电商、社交、电影、音乐、学术引用和知识图谱等6个领域，设计了8个评估任务（包括节点分类、链接预测、图分类等），评估了16个编码器和24个基线模型。OpenMAG的一个重要贡献是提供了统一的评估框架，支持不同编码器组合和融合策略的系统比较。

OpenMAG的实验进一步验证了MAGB的发现：多模态信息在不同领域的贡献存在差异，且简单的拼接融合策略并非最优选择。然而，OpenMAG在模态类型上仍局限于文本和图像双模态，同样不包含公式模态。此外，虽然OpenMAG声称覆盖了学术引用领域，但其学术数据集（如ogbn-arxiv）实际上并不包含原始的多模态属性，仅使用了OGB提供的低维文本向量。

\subsection{MM-GRAPH：模态消融基准}

Zhu等人\cite{zhu2025mmgraph}提出了MM-GRAPH，专注于模态消融实验的系统评估。MM-GRAPH包含5个多模态图数据集，支持7种模态组合的消融实验：仅图结构（G）、仅文本（T）、仅视觉（V）、图+文本（GT）、图+视觉（GV）、文本+视觉（TV）和全部模态（GTV）。这种设计使得研究者能够量化各模态的独立贡献和交互效应。

MM-GRAPH的消融实验揭示了一个重要现象：在大多数数据集上，文本+视觉的组合优于任何单一模态，但加入图结构信息后的三模态组合（GTV）并不总是优于双模态组合（TV）。这一发现表明，模态之间的交互效应并非简单的叠加关系，需要设计更精细的融合策略。然而，MM-GRAPH的领域同样局限于电商和社交领域，模态类型仅为文本和视觉两种，未涵盖公式等其他模态。

\subsection{现有MAG基准的共性问题}

综合分析MAGB\cite{yan2025magb}、OpenMAG\cite{wan2026openmag}和MM-GRAPH\cite{zhu2025mmgraph}三大MAG基准，可以发现它们存在以下共性局限：

\textbf{（1）模态类型单一。}现有MAG基准均只包含文本和图像两种模态，未涵盖公式这一学术论文特有的模态。公式承载了大量的结构化数学信息，在学术文献的分类和理解中具有独特价值。缺少公式模态使得现有MAG基准无法评估公式在学术图学习任务中的贡献。

\textbf{（2）领域覆盖有限。}现有MAG基准的数据集主要来源于电商和社交领域。学术论文作为一个天然的多模态属性图场景，其节点（论文）具有全文文本、实验图表、数学公式和表格等多种模态属性，其边（引用关系）具有明确的学术语义。电商和社交领域的发现（如"社交偏好视觉、电商偏好文本"）是否适用于学术领域，目前尚无系统的研究。

\textbf{（3）图结构差异。}电商和社交领域的图结构多为用户-商品交互图或社交关系图，其边的密度和分布特征与学术论文的引用网络有本质差异。引用网络通常具有有向性、稀疏性和幂律度分布等特征，这些特征对GNN消息传递的效果有重要影响。

本文提出的AMMAG正是为了填补上述空白：首次将公式模态引入MAG基准，将领域覆盖扩展到学术论文领域，并基于真实的引用网络结构构建图数据。

\section{多模态学术文档理解}

学术文档的多模态理解是自然语言处理和文档智能的重要研究方向。学术论文作为一种特殊的多模态文档，同时包含文本、图表、公式和表格等多种模态信息，其理解需要综合运用计算机视觉、自然语言处理和信息提取等多种技术。

\subsection{DocGenome：大规模科学文档基准}

Xia等人\cite{xia2024docgenome}提出的DocGenome是目前规模最大的科学文档多模态基准。DocGenome包含约50万篇科学文档，涵盖文本、图像、数学公式和页面布局四种模态，并提供了13种细粒度的布局属性标注和完整的LaTeX源代码。DocGenome覆盖了153个学科领域，从计算机科学到物理学、数学、生物学等，为训练和评估多模态大语言模型提供了丰富的数据基础。

DocGenome在文档布局分析方面采用了层次化的标注体系，将页面元素分为标题、正文、图片、图片标题、表格、表格标题、行内公式、行间公式、脚注、页眉、页脚等13种类型。这种细粒度的标注使得模型能够学习不同页面元素之间的语义关系。DocGenome的构建流程包括PDF解析、布局检测、内容提取和质量控制四个阶段，其中布局检测使用了DocLayout-YOLO等先进的文档布局分析模型。

DocGenome与AMMAG的关系体现在以下方面：一方面，DocGenome证明了学术论文确实包含丰富的多模态信息（特别是公式模态），这与本文的研究动机一致；另一方面，DocGenome专注于单文档级别的多模态理解任务（如文档分类、信息提取、视觉问答等），未构建文档间的引用网络图结构。AMMAG则将多模态属性与图结构相结合，首次在引用网络图上评估多模态图学习方法的性能。

\subsection{多模态科学问答与引用生成}

Li等人\cite{li2024m3sciqa}提出了M3SciQA，一个多文档多模态科学问答基准。M3SciQA要求模型同时理解多篇科学论文中的文本、图表和公式信息来回答复杂问题，强调跨文档的推理能力。M3SciQA包含多种问题类型，包括数值推理、图表解读、跨文档比较等，能够全面评估模型在科学文档理解方面的能力。M3SciQA的实验表明，当前最先进的多模态大语言模型在涉及公式理解和跨文档推理的问题上表现仍然不佳，说明学术文档的多模态理解仍是一个具有挑战性的任务。

Hu等人\cite{hu2025mcitebench}提出了MCiteBench，一个多模态引用生成基准。MCiteBench要求模型根据论文的文本和图表信息生成带有引用标记的文本，评估模型在学术写作辅助方面的能力。MCiteBench关注的是引用文本的生成质量，而非引用网络的结构特征。

Roberts等人\cite{roberts2024scifibench}提出了SciFIBench，专注于科学图表理解任务的评估。SciFIBench包含多种类型的科学图表（折线图、柱状图、散点图、示意图等），评估模型解读图表内容、提取数据趋势和进行科学推理的能力。SciFIBench的实验表明，即使是最先进的多模态模型在解读复杂科学图表时仍有较大提升空间。

\subsection{文档理解的技术基础}

多模态学术文档理解的技术基础包括文档布局分析、光学字符识别（OCR）和公式识别等。在文档布局分析方面，DocLayout-YOLO\cite{zhao2024doclayoutyolo}通过多样化的合成数据和全局-局部自适应感知机制，在科学文档的版面检测上取得了优异的性能。LayoutLMv3\cite{huang2022layoutlmv3}和DiT\cite{li2022dit}则通过预训练策略学习了文档图像中文本、布局和视觉信息的联合表示。Blecher等人\cite{blecher2024nougat}提出的Nougat模型能够将学术PDF文档转换为结构化的Markdown文本，保留了公式、表格和图片的语义信息。

在公式识别方面，UniMERNet\cite{wang2024unimernet}提出了一个通用的数学公式识别网络，能够将文档中的数学公式图像转换为LaTeX代码，在真实场景的公式识别任务上表现优异。在OCR方面，PaddleOCR\cite{paddleocr}是一个开源的多语言OCR工具包，支持多种文字识别场景。本文的自动化构建系统正是基于DocLayout-YOLO、UniMERNet和PaddleOCR等技术实现了学术论文的多模态信息提取。

\subsection{本节小结}

本节综述的多模态学术文档理解工作在文档级别的多模态信息提取和理解方面取得了重要进展。然而，这些工作关注的是单文档或少量文档的理解任务（如问答、引用生成、图表解读），均未将大量论文组织为图结构进行图学习任务的评估。学术论文间的引用关系构成了一种天然的图结构，利用这一图结构进行消息传递可能有助于更好地理解论文的语义。AMMAG正是将学术论文的多模态属性嵌入引用网络图结构中，首次在学术多模态属性图上开展系统评估。

\section{学术文献图构建与自动化采集}

学术文献图的构建涉及PDF文档解析、元数据提取、引用网络构建等多个环节，是实现学术多模态属性图数据集的基础。本节综述相关工具、数据集和自动化采集系统。

\subsection{学术PDF解析工具}

学术PDF文档的解析是构建学术文献图的第一步。Lopez\cite{lopez2009grobid}开发的GROBID（GeneRation Of BIbliographic Data）是目前使用最广泛的学术PDF结构化解析工具。GROBID基于条件随机场（CRF）和深度学习模型，能够从PDF文档中提取标题、作者、摘要、参考文献等结构化信息。GROBID采用模块化的架构设计，包含头部提取、参考文献解析、全文提取等多个模块，可以通过REST API进行调用。然而，GROBID是一种半自动工具，不支持图像和公式等多模态信息的提取，也无法自动构建完整的引用网络图。

Tkaczyk等人\cite{tkaczyk2015cermine}提出了CERMINE，一个基于规则和机器学习的学术PDF内容提取系统。CERMINE能够自动提取文档的元数据（标题、作者、摘要）、正文内容和参考文献，其准确率在多个标准数据集上与GROBID相当。ScienceParse是Allen人工智能研究所（AI2）开发的PDF解析工具，基于深度学习模型提取论文的结构化信息。与GROBID相比，ScienceParse更加轻量化，但在复杂文档格式上的处理能力稍弱。

这些传统PDF解析工具的共同局限在于：只能提取文本类型的元数据和内容，不支持图像、公式等多模态信息的提取；引用关系的提取依赖参考文献列表的解析，无法自动与论文数据库进行匹配构建完整的引用网络图；多为半自动工具，需要人工干预和配置，难以实现大规模的端到端自动化。

\subsection{大规模学术数据集}

大规模学术数据集为学术文献图的研究提供了重要的数据基础。Lo等人\cite{lo2020s2orc}提出的S2ORC（Semantic Scholar Open Research Corpus）是目前最大的开放学术语料之一，包含约4100万篇论文、约13亿条引用关系和约380万篇全文文本。S2ORC为每篇论文提供了结构化的元数据（标题、作者、摘要、发表年份、引用列表等），并将论文间的引用关系组织为引用网络图。然而，S2ORC主要提供文本类型的元数据，未提供论文的原始多模态内容（全文、图像、公式等）。

Tang等人\cite{tang2008arnetminer}开发的AMiner是一个学术搜索引擎和挖掘平台，提供了超过2亿篇论文的元数据搜索服务。AMiner构建了作者-论文-会议-机构的异构学术网络，支持学术关系挖掘和专家发现等任务。Priem等人\cite{priem2022openalex}提出的OpenAlex是一个开源的学术元数据仓库，提供了超过2.5亿篇学术作品的元数据和引用关系，覆盖论文、数据集、软件等多种学术成果类型。Clement等人\cite{clement2019arxiv}对arXiv数据集进行了系统分析，指出arXiv包含超过150万篇论文、约670万条引用边和110亿词的全文语料，是研究学术文献的重要数据源。

然而，上述数据集主要提供论文的元数据信息（标题、摘要、作者、引用关系等），未提供论文的原始多模态内容。研究者在使用这些数据集进行图学习研究时，只能使用低维的文本向量或词袋表示作为节点特征，丢失了大量有价值的多模态信息。

\subsection{自动化数据采集与多智能体系统}

随着大语言模型（LLM）技术的发展，基于多智能体系统的自动化数据采集成为可能。Ma等人\cite{ma2025autodata}提出的AutoData是一个基于多智能体系统的开放网络数据采集框架。AutoData利用多个大语言模型智能体协同完成数据的搜索、提取、验证和整合，实现了高度自动化的数据采集流程。AutoData的核心设计包括：搜索智能体负责从网络中发现相关数据源；提取智能体负责从网页中提取结构化数据；验证智能体负责检查数据的完整性和一致性；整合智能体负责将不同来源的数据合并为统一格式。然而，AutoData是一个通用框架，未针对学术多模态属性图的自动化构建进行专门设计。

Guo等人\cite{guo2024llm_agent_survey}系统综述了基于大语言模型的多智能体系统研究进展，将多智能体系统分为协作型、对抗型和层次型三种类型。在数据采集和处理的场景中，协作型多智能体系统最为适用——多个专用智能体各司其职，通过编排器协调执行顺序和数据流。本文的系统设计正是采用了这种协作型多智能体架构。

在图学习框架方面，Fey和Lenssen\cite{fey2019pyg}开发了PyTorch Geometric（PyG），提供了高效的GNN模型实现和数据处理工具。Wang等人\cite{wang2019dgl}开发了Deep Graph Library（DGL），支持多种深度学习框架后端。这些框架为图学习实验的实现提供了技术基础，本文的Benchmark实验即基于PyG框架实现。

\subsection{本节小结}

综合以上分析，现有学术文献图构建工具和数据采集系统存在以下不足：传统PDF解析工具（如GROBID）不支持多模态信息提取和引用网络自动构建；大规模学术数据集（如S2ORC）仅提供元数据，不包含原始多模态内容；通用多智能体框架（如AutoData）未针对学术MAG的自动化构建设计端到端流水线。本文提出的基于多智能体的自动化构建系统正是为了解决这些问题，实现了从arXiv元数据到标准化MAG数据集的端到端自动化流程。

\section{本章小结}

本章从四个方面系统综述了与本文研究密切相关的工作。在图神经网络与图基准方面，GCN\cite{kipf2017gcn}、GAT\cite{velickovic2018gat}、GraphSAGE\cite{hamilton2017graphsage}和GIN\cite{xu2019gin}等模型为图学习提供了坚实的方法基础，OGB\cite{hu2020ogb}等基准为评估提供了标准化框架，但现有基准在节点属性的丰富度上存在不足。在多模态属性图基准方面，MAGB\cite{yan2025magb}、OpenMAG\cite{wan2026openmag}和MM-GRAPH\cite{zhu2025mmgraph}开创了MAG学习范式，但均不包含公式模态且未覆盖学术论文领域。在多模态学术文档理解方面，DocGenome\cite{xia2024docgenome}等基准证明了学术论文包含丰富的多模态信息，但未将论文组织为图结构进行图学习评估。在自动化采集方面，GROBID\cite{lopez2009grobid}等工具和AutoData\cite{ma2025autodata}等系统提供了技术基础，但未针对学术MAG设计端到端流水线。

\begin{table}[htbp]
  \centering
  \caption{现有工作特性综合对比}
  \label{tab:related_work_final}
  \begin{tabular}{p{2.8cm}ccccccccc}
    \toprule
    特性 & \rotatebox{70}{OGB} & \rotatebox{70}{CS-TAG} & \rotatebox{70}{DocGenome} & \rotatebox{70}{MAGB} & \rotatebox{70}{OpenMAG} & \rotatebox{70}{MM-GRAPH} & \rotatebox{70}{S2ORC} & \rotatebox{70}{AutoData} & \rotatebox{70}{\textbf{AMMAG}} \\
    \midrule
    学术论文领域 & \checkmark & \checkmark & \checkmark & & & & \checkmark & & \textbf{\checkmark} \\
    节点多模态属性 & & & \checkmark & \checkmark & \checkmark & \checkmark & & & \textbf{\checkmark(3种)} \\
    图结构(引用网络) & \checkmark & \checkmark & & \checkmark & \checkmark & \checkmark & \checkmark & & \textbf{\checkmark} \\
    公式模态 & & & \checkmark & & & & & & \textbf{\checkmark} \\
    模态消融实验 & & & & \checkmark & \checkmark & \checkmark & & & \textbf{\checkmark(含公式)} \\
    自动化采集系统 & & & & & & & & \checkmark & \textbf{\checkmark} \\
    \bottomrule
  \end{tabular}
\end{table}

如表\ref{tab:related_work_final}所示，通过以上分析可以得出结论：目前没有任何一个工作同时满足学术论文领域、多模态节点属性（含公式模态）、引用网络图结构和自动化构建系统这四个条件。AMMAG正是填补这一空白的首次尝试，它继承了OGB和CS-TAG的学术领域和图结构优势，借鉴了DocGenome的多模态信息提取技术，融合了MAGB的MAG学习范式，并通过基于多智能体的自动化构建系统实现了端到端的数据集生成。
